\section{Four genus-one factors over the root orbifold}\label{sec:quadrics}

We construct the manifold in Proposition~\ref{prop:transfer}. No specialness assumption is used in the construction itself.

\begin{proposition}\label{prop:eightfold}
Let $X$ be a connected smooth projective fourfold and $D$ a nonempty smooth connected reduced divisor. There is a connected smooth projective eightfold $Y$ with a map to $\sqrt[2]{(X,D)}$ whose coarse map $g:Y\to X$ satisfies all the hypotheses of Proposition~\ref{prop:transfer}. Its smooth general fibre is a product of four genus-one curves.
\end{proposition}

The parameter choice has two roles. Four independent factors eliminate stabilizers over the three-dimensional divisor, while a separate rank estimate supplies a submersive point over the general point of every base divisor. Both are needed in Proposition~\ref{prop:transfer}; generic smoothness alone would not justify the descent of all differential tensors.

Choose a line bundle $M$ such that $M$ and $M(-D)$ are globally generated. This is possible by taking a sufficiently high power of an ample line bundle. On $\cX$, with $T$ as in \eqref{eq:root-line}, put
\[
 V=\OO_{\cX}^{\oplus2}\oplus T^{\oplus2},\qquad
 \mathcal P=\underbrace{\PP_{\cX}(V)\times_{\cX}\cdots\times_{\cX}\PP_{\cX}(V)}_{4\text{ factors}}.
\]
All four factors are over the \emph{same} root orbifold. In a root chart there is one diagonal sign action on their product.

In factor $j$, use vector coordinates $z_{j1},\ldots,z_{j4}$ and choose two diagonal quadratic equations
\begin{equation}\label{eq:quadrics}
 q_{j\nu}=\sum_{i=1}^4a_{j\nu i}z_{ji}^2=0,
 \qquad j=1,\ldots,4,\quad \nu=1,2,
\end{equation}
where, independently for all $j,\nu,i$,
\[
 a_{j\nu i}\in
 \begin{cases}
 H^0(X,M),&i=1,2,\\
 H^0(X,M(-D)),&i=3,4.
 \end{cases}
\]
The squares of the last two coordinates take values in $T^2=\pi^*\OO_X(D)$, so each equation is a section of the appropriate $\OO(2)\otimes\pi^*M$. Let $\mathsf A$ denote the affine space of all these coefficients, and let $\mathcal S_a\subset\mathcal P$ be the simultaneous zero locus for $a\in\mathsf A$.

In compatible frames at $x\in X$, write the two rows in factor $j$ as a $2\times4$ matrix $C_j(x)$. The ranks of its specified column submatrices are independent of the choice of frames. Evaluation onto the four matrices is surjective at every $x$, by global generation and independence of the coefficient choices. This includes $x\in D$: a coefficient in $M(-D)$ is evaluated in the fibre of that line, and is paired with a square in $\OO_X(D)$. Its numerical value is not required to vanish on $D$.

\begin{lemma}[Simultaneous parameter choice]\label{lem:parameters}
There is a choice $a\in\mathsf A$ with the following properties.
\begin{enumerate}
\item\label{item:total-smooth} The zero locus $\mathcal S_a$ is smooth in its root charts.
\item\label{item:free} It has no nontrivial stabilizers.
\item\label{item:triples} There is a closed subset $Z\subset X$ of codimension at least two such that every matrix obtained from any $C_j(x)$ by deleting one column has rank two for $x\notin Z$.
\item\label{item:pairs} On a nonempty open $B\subset X\setminus D$, every $2\times2$ minor of every $C_j(x)$ is nonzero.
\end{enumerate}
\end{lemma}

\begin{proof}
Take a finite algebraic trivializing cover of $X$ for the bundles involved, with smooth root charts
\[
 W=\{(x,w):w^2=t(x)\}\longrightarrow U,
\]
where $t$ is the local divisor equation. These charts are smooth since $D$ is smooth and reduced; away from $D$ they are \emph{\'{e}tale}. Their products with the four projective spaces are smooth. At any point of such a product, each equation in \eqref{eq:quadrics} has nonzero evaluation as a linear functional on its own parameter block: at least one coordinate in that projective factor is nonzero. The eight evaluations are independent. Thus the universal incidence over the chart product is the kernel of a surjective vector-bundle map from the parameter bundle to the eight equation lines, and is smooth.

Generic smoothness in characteristic zero makes its fibres over a dense parameter open smooth. Indeed, each component of the critical locus has image of dimension smaller than $\dim\mathsf A$, since a dominating component would have generically surjective differential on its smooth locus, contradicting criticality. Images of nondominating components may likewise be discarded. There are finitely many algebraic charts, so one common dense open of parameters proves~\ref{item:total-smooth}.

Only points over $D$ can have nontrivial stabilizer. There the involution on a factor sends
\[
 [z_1:z_2:z_3:z_4]\longmapsto[z_1:z_2:-z_3:-z_4].
\]
Its fixed locus is the union of the two eigenspace lines. If a fixed point satisfies both equations, either the first pair of columns or the last pair of columns of $C_j(x)$ is dependent. Each of these determinant conditions is a proper hypersurface in matrix space. A fixed point in the fourfold product therefore requires one such condition in \emph{each} of the four independent matrices. At fixed $x\in D$ its parameter codimension is four. The bad incidence over $D$ has dimension at most
\begin{equation}\label{eq:fixed-dimension}
 \dim\mathsf A+\dim D-4=\dim\mathsf A-1.
\end{equation}
General parameters avoid its image. This excludes stabilizers above every point of $D$, including degenerate fibres, and proves~\ref{item:free}.

For one specified three-column submatrix, rank at most one has codimension two in the space of $2\times3$ matrices. Surjective evaluation implies that its incidence over $X$ has dimension at most $\dim\mathsf A+2$. For a general parameter its bad base locus has dimension at most two, by the fibre-dimension theorem. The finitely many choices of factor and deleted column can be treated simultaneously. Their union is closed and gives~\ref{item:triples}.

Finally choose a point $x_0\in X\setminus D$. By evaluation surjectivity, all pair minors can be made nonzero at $x_0$; for example the columns $(1,\lambda_i)^{\mathsf t}$ with distinct $\lambda_i$ have this property. This is a nonempty Zariski-open parameter condition and can be intersected with the preceding dense opens. Nonvanishing then holds on an open neighbourhood $B$ of $x_0$, proving~\ref{item:pairs}.
\end{proof}

Fix such a parameter and write $\mathcal S=\mathcal S_a$. Freeness and smoothness mean that its chart quotients glue to an ordinary smooth complex space $S$ with a map $S\to\cX$. In particular the quotient of each source chart is covered \emph{unramifiedly} by that chart; its map to the base root chart gives the local lifts on $S$.

\begin{lemma}[Projectivity]\label{lem:projectivity}
The space $S$ is projective.
\end{lemma}

\begin{proof}
Square the four coordinates in each factor. The identification $T^2=\pi^*\OO_X(D)$ gives a globally defined invariant map to the ordinary projective variety
\begin{equation}\label{eq:square-target}
 Q=\underbrace{\PP_X(E_0)\times_X\cdots\times_X\PP_X(E_0)}_{4\text{ factors}},
 \qquad E_0=\OO_X^{\oplus2}\oplus\OO_X(D)^{\oplus2}.
\end{equation}
Its map on a root chart is the finite root map $W\to U$ together with the finite projective coordinate-square maps. Restriction to the closed zero locus remains finite. Passing to the finite-group quotient preserves finiteness: over an affine target this takes an invariant subalgebra of a finite algebra, which is again finite. These invariant algebras glue under the bundle transitions, since the coordinate-square construction is invariant and global. Hence $S\to Q$ is a finite analytic map.

For completeness, a finite analytic space over a complex projective variety algebraizes by projective GAGA \cite{GAGA}. Its finite coherent direct-image algebra, including multiplication and unit, algebraizes, and relative spectrum recovers the space after analytification. A finite algebraic morphism is projective \cite[Tag~01WG]{Stacks}. Thus $S$ is projective and in particular compact.
\end{proof}

\begin{lemma}[Fibres and submersions over base divisors]\label{lem:fibres-witnesses}
The map $S\to X$ is smooth over $B$ with connected fibres that are products of four genus-one curves. Above every root-chart point with image in $X\setminus Z$, it has a point at which the lifted map to that chart is a submersion.
\end{lemma}

\begin{proof}
At $x\in B$, consider one matrix $C=C_j(x)$. It has rank two. Coordinate squaring on $\PP^3$ is finite and surjective; the common quadric zero set is its inverse image of $\PP(\ker C)\simeq\PP^1$. Thus it is nonempty and has dimension one. At a solution, put $b_i=z_i^2$. A nonzero vector $b\in\ker C$ cannot have support one or two, since every pair of columns is independent. The vertical derivative
\begin{equation}\label{eq:vertical-jacobian}
 2C\operatorname{diag}(z_1,z_2,z_3,z_4)
\end{equation}
therefore has rank two. It annihilates the radial direction at a solution, so it still has rank two on the projective tangent space. This proves both smoothness of the curve and submersivity of the family over $B$.

The curve is a smooth complete intersection of two quadrics in $\PP^3$. Its Koszul resolution is
\[
 0\longrightarrow\OO_{\PP^3}(-4)
 \longrightarrow\OO_{\PP^3}(-2)^{\oplus2}
 \longrightarrow\OO_{\PP^3}
 \longrightarrow\OO_C\longrightarrow0.
\]
The intermediate cohomology vanishings for line bundles on $\PP^3$ give $H^0(C,\OO_C)=\CC$, hence connectedness \cite[Tag~01XS]{Stacks}. Adjunction gives $K_C=\OO_C$. Thus $C$ is a genus-one curve and its cotangent line is trivial. Taking four factors proves the first assertion and triviality of the cotangent bundle of each product fibre.

Now let $x\notin Z$, possibly on $D$. In each matrix $C_j(x)$, the three columns remaining after any deletion span $\CC^2$. Its kernel is consequently contained in no coordinate hyperplane: one can prescribe a chosen coordinate to be one and solve for the other three. A vector space over $\CC$ is not the union of finitely many proper linear subspaces, so its kernel contains a vector $b$ with all four coordinates nonzero. Choose square roots $z_i$ of its entries. Formula~\eqref{eq:vertical-jacobian} then has rank two. Choosing such a point in each factor makes the full vertical derivative of the eight equations surjective. The implicit function theorem gives a submersion from their zero set to the base root chart at this tuple. The argument works for either chosen chart point over $x$ and in particular for $w=0$.
\end{proof}

\begin{proof}[Proof of Proposition~\ref{prop:eightfold}]
The smooth projective $S$ has disjoint connected components. Its fibres over $B$ are connected and nonempty, so there is exactly one component $Y$ dominating $X$. More explicitly, the proper image of a dominating component is all of $X$; two such components would disconnect the fibres over $B$. Every submersive point constructed in Lemma~\ref{lem:fibres-witnesses} lies on a dominating component, since its local image has full dimension. Therefore all these points lie on $Y$.

The variety $Y$ is connected smooth projective, of dimension $4+4=8$. Its map $g:Y\to X$ is surjective and has the asserted fibres over $B$. Its map to $\cX$ gives
\[
 g^*D=2f^*\cD.
\]
The pullback of the root-divisor section is generically nonzero, so $f^*\cD$ is an effective integral Cartier divisor. Thus every divisorial multiplicity over $D$ is even. Since $\codim_X Z\ge2$, every base prime divisor has general point outside $Z$. Lemma~\ref{lem:fibres-witnesses} supplies the ordinary or root-chart submersions there. All four conditions of Proposition~\ref{prop:transfer} follow.
\end{proof}

\begin{proof}[Proof of Theorem~\ref{thm:main}]
Apply Proposition~\ref{prop:eightfold} to $(X,D)$. Specialness of $\cX$ and Proposition~\ref{prop:transfer} make the resulting eightfold $Y$ special and give $\pi_1(Y)\twoheadrightarrow G$. Assumption~\ref{ass:manifold}, in dimension eight, makes $\pi_1(Y)$ virtually abelian. Its quotient $G$ is virtually abelian as claimed.
\end{proof}
